\section{Methodology}
\label{sec:experiments}

This section details the experimental setup designed to validate the vulnerability of modern time-series forecasters to both black-box and white-box sponge attacks.

\subsection{Dataset and Preprocessing}
We utilize the \textbf{Wind Power Generation Data} dataset acquired from Kaggle, which provides high-resolution time-series data for renewable energy forecasting.
\begin{itemize}
    \item \textbf{Features:} The dataset includes critical meteorological and operational features such as Wind Power (MW), Wind Speed (m/s), Wind Direction ($^\circ$), and ambient temperature.
    \item \textbf{Normalization:} To ensure stable gradient descent during the attack generation, all continuous features were scaled using Z-score normalization.
    \item \textbf{Input Window:} We structured the data into sliding windows of length $T=64$ hours to predict the subsequent $T=10$ hours, enabling the model to capture diurnal cycles and weather fronts.
\end{itemize}

\subsection{Target Architectures}
We contest three distinct architectures to test our hypothesis regarding Static versus Dynamic computation graphs.
\begin{itemize}
    \item \textbf{DeepAR:} A standard 2-layer LSTM (Hidden Size: 128) representing static graphs.
    \item \textbf{Chronos-v2:} A pretrained T5-based foundation model representing quantized, deterministic graphs.
    \item \textbf{ACT-LSTM:} A custom implementation of Adaptive Computation Time, representing dynamic graphs with vulnerable control flow.
\end{itemize}

\subsection{Attack Implementation}
To evaluate the full spectrum of risk defined in our threat model, we implement two distinct attack vectors.

\subsubsection{Baseline: Genetic Algorithm (Black-Box)}

\begin{figure}[t]
\centering
\begin{tikzpicture}[
    node distance=1.2cm and 1.5cm,
    every node/.style={font=\sffamily\footnotesize, align=center},
    % Styles
    block/.style={rectangle, draw=black!80, thick, rounded corners, minimum width=2cm, minimum height=1cm, fill=white, drop shadow},
    process/.style={rectangle, draw=blue!80, thick, rounded corners, minimum width=2.5cm, minimum height=1.2cm, fill=blue!5},
    cycle/.style={circle, draw=green!60!black, thick, fill=green!5, minimum size=2.5cm},
    arrow/.style={-Latex, thick, draw=black!70}
]

    % --- 1. The Target (Opaque) ---
    \node[block, fill=black!80, text=white] (model) {\textbf{Target Model} \\ (Black Box) \\ \tiny Unknown Weights $\theta$};

    % --- 2. The Interaction ---
    % Input Population
    \node[block, left=of model, label=above:\textbf{Population $P_t$}] (pop) {
        \faChartLine \ \faChartLine \\[-0.3em] \tiny Candidate Inputs
    };

    % Output Feedback
    \node[right=of model, text width=2cm] (feedback) {Latency $L(x)$ \\ \tiny (Fitness Score)};

    % --- 3. The Genetic Engine (Evolutionary Cycle) ---
    \node[process, below=1.5cm of model] (ga) {\textbf{Genetic Algorithm} \\ \tiny Selection $\rightarrow$ Crossover $\rightarrow$ Mutation};

    % --- CONNECTIONS ---

    % Forward Pass (Query)
    \draw[arrow] (pop) -- node[above, font=\tiny] {Query} (model);
    \draw[arrow] (model) -- (feedback);

    % Feedback Loop
    \draw[arrow] (feedback.south) |- (ga.east);
    \draw[arrow] (ga.west) -| (pop.south);

    % Annotations
    \node[below=0.1cm of ga, font=\tiny, color=black!60] {Repeat $N$ Generations};

\end{tikzpicture}
\caption{\textbf{Black-Box Threat Model (Genetic Algorithm).} The adversary has no access to model internals. They maintain a population of input waveforms, query the black-box model to measure execution time, and evolve the population to maximize latency.}
\label{fig:blackbox_attack}
\end{figure}

As described in previous work \cite{shumailov2021sponge}, we utilize a Genetic Algorithm (PyGAD) to evolve adversarial inputs. The fitness function is defined solely by the wall-clock inference time: $F(x) = -\text{Latency}(x)$. This assumes $\mathcal{K}_{BB}$.

\subsubsection{Proposed: XAI-Guided Saliency Attack (White-Box)}
\begin{figure}[t]
\centering
\begin{tikzpicture}[
    node distance=1.2cm and 1.5cm,
    every node/.style={font=\sffamily\footnotesize, align=center},
    % Styles
    block/.style={rectangle, draw=black!80, thick, rounded corners, minimum width=2cm, minimum height=1cm, fill=white, drop shadow},
    whitebox/.style={rectangle, draw=orange!80, thick, rounded corners, minimum width=2.5cm, minimum height=1.5cm, fill=white},
    arrow/.style={-Latex, thick, draw=black!70},
    gradient/.style={draw=red!80, thick, dashed, color=red}
]

    % --- 1. The Target (Transparent) ---
    \node[whitebox] (model) {};
    % Internals visible
    \node[font=\tiny, anchor=north] at (model.north) {\textbf{Target Model} (White Box)};
    \node[draw, circle, inner sep=1pt, minimum size=0.5cm, fill=orange!10] at (model.center) (loop) {$\circlearrowright$};
    \node[below=0.1cm of loop, font=\tiny] {Ponder Loop $h(s_t)$};

    % --- 2. The Flow ---
    \node[block, left=of model] (input) {$X + \delta$};
    \node[right=of model, text width=2.5cm] (loss) {\textbf{Ponder Loss} \\ $\mathcal{L} = \sum p_t$};

    % --- 3. The Adversary (Gradient Step) ---
    \node[block, below=1.5cm of model, fill=red!5, draw=red!50] (attacker) {
        \textbf{PGD Update} \\
        \tiny $x \leftarrow x + \alpha \cdot \text{sign}(\nabla_x \mathcal{L})$
    };

    % --- CONNECTIONS ---

    % Forward Pass
    \draw[arrow] (input) -- (model);
    \draw[arrow] (model) -- (loss);

    % Gradient Backprop (The Attack)
    \draw[arrow, gradient] (loss.south) |- node[right, font=\tiny, pos=0.2] {Backprop $\nabla_\theta$} (attacker.east);
    \draw[arrow, gradient] (attacker.west) -| node[left, font=\tiny, pos=0.2] {Perturb $\delta$} (input.south);

    % Visuals
    \node[above=0.1cm of input] {\faSyringe}; % Injection icon

\end{tikzpicture}
\caption{\textbf{White-Box Threat Model (Saliency-Guided PGD).} The adversary inspects the internal halting probabilities (Ponder Loss). By computing the gradient of computation time with respect to the input ($\nabla_x \mathcal{L}$), they iteratively optimize the perturbation $\delta$ to maximize the loop duration.}
\label{fig:whitebox_attack}
\end{figure}

For the $\mathcal{K}_{WB}$ scenario, we propose a novel \textit{Saliency-Guided Sponge Attack}. Instead of maximizing classification error, we define a "Ponder Loss" $\mathcal{L}_{ponder}$ representing the total accumulated halting probability:
\begin{equation}
    \mathcal{L}_{ponder}(x) = \sum_{t=1}^{T} \sum_{k=1}^{M} h(s_{t,k})
\end{equation}
We compute the gradient $\nabla_x \mathcal{L}_{ponder}$ to generate a Saliency Map. This map identifies the specific time-steps that trigger the halting unit. We then apply Projected Gradient Descent (PGD) masked by this saliency map, effectively "pushing" the sensitive time-steps towards the decision boundary where the model hesitates the most.

\subsection{Formal Algorithm: Saliency-Guided PGD}
To ensure reproducibility, we formalize our White-Box attack procedure in Algorithm \ref{alg:sponge}. Unlike standard PGD which maximizes classification loss $\mathcal{L}_{CE}$, our algorithm maximizes the Ponder Cost $\mathcal{L}_{ponder}$ (Eq. 5) while respecting the physical constraints of the wind sensor data (e.g., wind speed cannot be negative).

\begin{algorithm}[h]
\caption{Saliency-Guided Sponge Attack (PGD)}
\label{alg:sponge}
\begin{algorithmic}[1]
\Require Model $f_\theta$, Input $x$, Step size $\alpha$, Iterations $K$, Max Perturbation $\epsilon$
\Ensure Sponge Example $x_{adv}$

\State $x_{adv} \leftarrow x$
\State $x_{orig} \leftarrow x$

\For{$k = 1$ to $K$}
    \State \textit{// Forward pass to get Halt Probabilities}
    \State $N, \{h_t\} \leftarrow f_\theta(x_{adv})$

    \State \textit{// Compute Ponder Loss (Sum of Halting Probs)}
    \State $\mathcal{L}_{ponder} \leftarrow \sum_{t} \sum_{step} h_t(step)$

    \State \textit{// Compute Gradient w.r.t Input}
    \State $\nabla_x \mathcal{L} \leftarrow \frac{\partial \mathcal{L}_{ponder}}{\partial x_{adv}}$

    \State \textit{// Compute Saliency Mask (Top 10\% features)}
    \State $M \leftarrow \mathbb{I}(|\nabla_x \mathcal{L}| > \text{percentile}(|\nabla_x \mathcal{L}|, 90))$

    \State \textit{// Update with Masked Gradient Ascent}
    \State $x_{adv} \leftarrow x_{adv} + \alpha \cdot \text{sign}(\nabla_x \mathcal{L}) \cdot M$

    \State \textit{// Project back to $\epsilon$-ball and valid range}
    \State $\delta \leftarrow \text{clip}(x_{adv} - x_{orig}, -\epsilon, \epsilon)$
    \State $x_{adv} \leftarrow \text{clip}(x_{orig} + \delta, 0, \text{max\_val})$
\EndFor

\State \Return $x_{adv}$
\end{algorithmic}
\end{algorithm}

The key innovation in Algorithm \ref{alg:sponge} is the masking step (Line 9). By filtering the gradient update $M$, we ensure that perturbations are applied \textit{only} to the time-steps that strongly influence the halting unit. This prevents the "high-frequency noise" look of standard adversarial examples, making our sponge inputs visually smoother and harder to detect.

\subsubsection{Bit-Flip Oracle Attack (Energy via Switching Activity)}
While the PGD and GA attacks increase energy consumption indirectly by prolonging computation time, we additionally target the \textit{dynamic power dissipation} directly. The Bit-Flip Oracle approximates the switching activity on the data bus during the first matrix multiplication of inference.

Given an input $x \in \mathbb{R}^{W \times F}$ and the first-layer weight matrix $W_1 \in \mathbb{R}^{m \times n}$, the oracle first tiles the flattened input to match the weight dimensions:
\begin{equation}
    x_{tiled} = \text{tile}(\text{flatten}(x), |W_1|)
\end{equation}
It then converts both tensors to their IEEE 754 binary representation and counts the bit positions that differ via XOR:
\begin{equation}
    \text{Flips}(x) = \sum_{i} \text{popcount}(\text{bits}(x_{tiled})_i \oplus \text{bits}(W_1)_i)
\end{equation}
Each bit difference corresponds to a transistor switching event during the multiply-accumulate operation, contributing to dynamic power dissipation ($P_{dyn} \propto \alpha C V^2 f$, where $\alpha$ is the switching activity factor).

Since the input is tiled to the full weight matrix size, the flip count scales proportionally with the number of parameters in the first layer. This means larger models inherently produce more switching activity per inference, even under benign inputs.

We optimize this metric using a Genetic Algorithm with fitness $F(x) = \text{Flips}(x) / \text{Flips}(x_{baseline})$, evolving inputs that maximize switching activity while respecting the same physical constraints as the latency and energy attacks.
